% Generated draft: requires mathematical and source review.
\documentclass[11pt]{article}
\usepackage[margin=25mm]{geometry}
\usepackage{fontspec}
\setmainfont{Latin Modern Roman}
\usepackage{amsmath,amssymb,mathtools,unicode-math}
\setmathfont{Latin Modern Math}
\usepackage{enumitem,needspace,xurl,hyperref}
\setlength{\parindent}{0pt}
\setlength{\parskip}{5pt}
\setlength{\emergencystretch}{3em}
\newcommand{\N}{\mathbb{N}}
\newcommand{\Z}{\mathbb{Z}}
\newcommand{\Q}{\mathbb{Q}}
\newcommand{\R}{\mathbb{R}}
\newcommand{\C}{\mathbb{C}}
\newcommand{\F}{\mathbb{F}}
\newcommand{\A}{\mathbb{A}}
\newcommand{\PP}{\mathbb{P}}
\newcommand{\E}{\mathbb{E}}
\newcommand{\eps}{\varepsilon}
\newcommand{\dd}{\,\mathrm{d}}
\begin{document}
\section*{UnsolvedMath ID: 701620}
\textbf{Source problem code:} catalogue entry
\subsection{Definitions and mathematical statement}
The supplied record does not define the class \(\mathscr{S}\) or specify a topology on its groups. A contraction group requires topological data, so the question cannot be formulated unconditionally from the supplied record alone.

For a topological group \(G\) with identity \(e_G\), use the convention
\[
\operatorname{con}_G(g)=\{x\in G: g^n x g^{-n}\longrightarrow e_G\text{ as }n\to\infty\},\qquad g\in G.
\]
The contraction group of \(g\) is called trivial when \(\operatorname{con}_G(g)=\{e_G\}\).

Subject to recovering the intended definition of \(\mathscr{S}\) and its topological hypotheses, the apparent existence question is:
\[
\text{Does there exist }G\in\mathscr{S}\text{ such that }
\forall g\in G,\quad \operatorname{con}_G(g)=\{e_G\}\,?
\]
This is a conditional reconstruction: the missing definition of \(\mathscr{S}\) must supply all restrictions on the admissible groups and their topologies. The supplied wording does not explicitly distinguish this existence question from a question about an arbitrarily specified member of \(\mathscr{S}\).
\subsection{Short English statement}
The apparent question asks whether some group in the intended class \(\mathscr{S}\) can have trivial contraction group for every element. The record leaves \(\mathscr{S}\) and its topological hypotheses undefined, so the formulation remains conditional.
\subsection{Sources}
\begin{itemize}
\item Supplied local UnsolvedMath registry record: permanent ID \(701620\), exact source problem code \texttt{catalogue entry}, title \emph{Kourovka Notebook Problem 19.72}.
\item The record attributes Problem 19.72 to C. Reid, proposed in 2018, and cites \emph{The Kourovka Notebook (Unsolved Problems in Group Theory)}, 19th issue (2018), with the 21st issue edition (2026), \texttt{arXiv:1401.0300v48}. The supplied collection description identifies E. I. Khukhro and V. D. Mazurov as editors.
\item The registry labels the problem \emph{open}. This is supplied source metadata, not an independently verified mathematical or literature assessment. The cited source was not independently consulted for this draft.
\end{itemize}
\end{document}
